\section{A227013: harmonic means between fourth powers}
\label{sec:quartic-harmonic}

For $n\ge1$, define the literal harmonic mean and its floor by
\[
 d_n=n^4-(n-1)^4,\qquad
 H_n=\sum_{k=(n-1)^4+1}^{n^4}\frac1k,\qquad
 M_n=\frac{d_n}{H_n},\qquad a_n=\lfloor M_n\rfloor.
\]
Kimberling's A227013 records a conjectural rational generating function,
closed form, and recurrence \cite{oeis227013}. We prove the exact generating
function, while explicitly correcting the index conventions of the other
formulas. This is one source family.

\begin{theorem}\label{thm:quartic-main}
Put
\[
 Q(n)=n^4-2n^3+\frac53n^2-\frac23n
      =\frac{n(n-1)(3n^2-3n+2)}3.
\]
Then $a_1=1$, $a_2=6$, and $a_n=\lfloor Q(n)\rfloor$ for every $n\ge3$.
The exact formal generating function, with coefficient $x^j$ equal to
$a_{j+1}$, is
\begin{equation}\label{eq:quartic-gf}
 \sum_{j\ge0}a_{j+1}x^j=
 \frac{-1-2x-22x^2-23x^3-20x^4-4x^5+2x^6-3x^7+x^8}
 {(x-1)^5(1+x+x^2)}.
\end{equation}
For every $n\ge10$,
\begin{equation}\label{eq:quartic-recurrence}
 a_n=4a_{n-1}-6a_{n-2}+5a_{n-3}-5a_{n-4}
       +6a_{n-5}-4a_{n-6}+a_{n-7}.
\end{equation}
\end{theorem}

\subsection{A finite rational moment bound}

Consider any consecutive block $a+1,\ldots,a+d$ with $a\ge0$ and
positive integer $d$. Set
\[
 c=a+\frac{d+1}2,\qquad r=\frac{d-1}2,\qquad t_j=j-r\quad(0\le j<d).
\]
The reflection $j\mapsto d-1-j$ negates $t_j$. Thus
\begin{equation}\label{eq:quartic-pairing}
 \sum_{j=0}^{d-1}\frac1{a+1+j}
 =\sum_{j=0}^{d-1} B(c,t_j),\qquad
 B(c,t)=\frac12\left(\frac1{c+t}+\frac1{c-t}\right).
\end{equation}
Here $c>r\ge|t_j|$, so all denominators are positive.
The exact identities
\begin{align}
 B(c,t)-\left(\frac1c+\frac{t^2}{c^3}+\frac{t^4}{c^5}+\frac{t^6}{c^7}\right)
 &=\frac{t^8}{c^7(c^2-t^2)}\ge0,\label{eq:quartic-lower-pair}\\
 \left(\frac1c+\frac{t^2}{c^3}+\frac{t^4}{c^3(c^2-r^2)}\right)-B(c,t)
 &=\frac{t^4(r^2-t^2)}{c^3(c^2-r^2)(c^2-t^2)}\ge0
 \label{eq:quartic-upper-pair}
\end{align}
hold for every $|t|\le r$. No infinite expansion or analytic remainder
estimate is involved.

The averaged central moments of $0,\ldots,d-1$ are
\begin{align*}
 v_2&=\frac1d\sum_{j<d}t_j^2=\frac{d^2-1}{12},\\
 v_4&=\frac1d\sum_{j<d}t_j^4=\frac{(d^2-1)(3d^2-7)}{240},\\
 v_6&=\frac1d\sum_{j<d}t_j^6=\frac{(d^2-1)(3d^4-18d^2+31)}{1344}.
\end{align*}
These identities follow by expanding the powers and summing the ordinary
powers through degree six. In the formal proof each ordinary power sum
is proved by induction on $d$, and the three displayed identities follow
by polynomial algebra. Summing \eqref{eq:quartic-lower-pair} and
\eqref{eq:quartic-upper-pair} gives the exact bounds
\begin{equation}\label{eq:quartic-LU}
 dL\le H\le dU,\qquad
 L=\frac1c+\frac{v_2}{c^3}+\frac{v_4}{c^5}+\frac{v_6}{c^7},\qquad
 U=\frac1c+\frac{v_2}{c^3}+\frac{v_4}{c^3(c^2-r^2)}.
\end{equation}

\subsection{Uniform separation of the harmonic mean from its floor}

For the quartic block, take $a=(n-1)^4$ and $d=d_n$. Then
\[
 c=(n^2-n+1)^2,\qquad r=(d_n-1)/2,\qquad
 c^2-r^2=n^4((n-1)^4+1).
\]
For $n=t+5$, direct clearing of denominators in \eqref{eq:quartic-LU}
gives the following two exact rational identities:
\begin{align}
 (Q(n)+2/3)L-1
 &=\frac{P(t)}{315(t^2+9t+21)^{14}},\label{eq:quartic-P}\\
 1-Q(n)U
 &=\frac{R(t)}{45(t+5)^2(t^2+9t+21)^6((t+4)^4+1)}.
 \label{eq:quartic-R}
\end{align}
All coefficients of $P$ and $R$, including their constant terms, are
positive. The complete coefficient lists appear in Table~\ref{tab:quartic-cert};
they define the polynomials by $P(t)=\sum_kP_kt^k$ and
$R(t)=\sum_kR_kt^k$. Consequently both right sides are strictly positive
for every $t\ge0$. The identities themselves are checked by polynomial
normalization in Lean, so the table is an exact algebraic certificate.

Since $Q(n)>0$, $d_n>0$, and $H_n>0$, equations
\eqref{eq:quartic-LU}--\eqref{eq:quartic-R} imply
\[
 Q(n)H_n\le Q(n)d_nU<d_n,
 \qquad
 (Q(n)+2/3)H_n\ge (Q(n)+2/3)d_nL>d_n.
\]
Therefore, for every $n\ge5$,
\begin{equation}\label{eq:quartic-sandwich}
 Q(n)<M_n<Q(n)+\frac23.
\end{equation}

For an integer $n$, the fractional part of $Q(n)$ is $1/3$ when
$n\equiv2\pmod3$ and is zero otherwise. This follows by substituting
$n=3k,3k+1,3k+2$ into its polynomial. Thus
\[
 Q(n)+2/3\le\lfloor Q(n)\rfloor+1.
\]
Together with \eqref{eq:quartic-sandwich}, this places $M_n$ in the
same unit interval as $Q(n)$ and proves $a_n=\lfloor Q(n)\rfloor$
for every $n\ge5$.

The remaining literal harmonic means satisfy the following finite rational
comparisons:
\[
\begin{array}{c|r|r|c}
n&d_n&a_n&\text{exact comparison for }H_n\\\hline
1&1&1&H_1=1\\
2&15&6&15/7<H_2\le15/6\\
3&65&40&65/41<H_3\le65/40\\
4&175&152&175/153<H_4\le175/152
\end{array}
\]
They can be checked by multiplying through by the common denominator
of the finitely many positive reciprocals. In Lean all four values
are explicit kernel computations on exact rational numbers. In particular
$\lfloor Q(3)\rfloor=40$ and $\lfloor Q(4)\rfloor=152$, completing
the floor formula at every claimed index.

\begin{table}[p]
\centering
\small
\caption{Complete positive-coefficient certificates for
\eqref{eq:quartic-P} and \eqref{eq:quartic-R}. A dash means coefficient zero.}
\label{tab:quartic-cert}
\begin{tabular}{r|r|r}
$k$&$P_k$&$R_k$\\\hline
24&28&---\\
23&3024&---\\
22&156572&---\\
21&5172156&---\\
20&122369371&---\\
19&2206424754&---\\
18&31496055869&---\\
17&364973853078&---\\
16&3492525189361&---\\
15&27925808790492&---\\
14&188059104987530&26\\
13&1071825627742992&1612\\
12&5182093657336537&46362\\
11&21253816999367298&819836\\
10&73776484843852949&9958385\\
9&215711119121472558&87898188\\
8&527200074613081231&581403224\\
7&1064965904311292064&2927810658\\
6&1749657582705865988&11280438711\\
5&2284392530671416108&33094892742\\
4&2290486550158893280&72791311130\\
3&1671930872839271052&116416933660\\
2&810572143783759280&128020606587\\
1&216472755572592948&86680806370\\
0&18186223121964135&27278343429\\
\end{tabular}
\end{table}

\subsection{The recurrence and exact generating function}

The sequence $b_n=\lfloor Q(n)\rfloor$ is a quartic polynomial plus
a correction of period three:
\[
 b_n=Q(n)-\frac13\mathbf1_{n\equiv2\ (\mathrm{mod}\ 3)}.
\]
If $E$ is the forward shift, the polynomial part is annihilated by
$(E-1)^5$ and the periodic part by $E^3-1$. Both are therefore
annihilated by
\[
 (E-1)^5(E^2+E+1)
 =E^7-4E^6+6E^5-5E^4+5E^3-6E^2+4E-1.
\]
This gives \eqref{eq:quartic-recurrence} whenever all its indices are
at least three, that is, $n\ge10$. The Lean proof expands the three
residue-class polynomials and checks this identity for arbitrary
quotient $k$, then transports it to the literal source sequence.

The first nine terms are
\[
 (a_1,\ldots,a_9)=(1,6,40,152,413,920,1792,3173,5232).
\]
Let $A(x)=\sum_{j\ge0}a_{j+1}x^j$ and
$D(x)=(1-x)^5(1+x+x^2)$.
The recurrence makes every coefficient of $D(x)A(x)$ from degree nine
onward zero; its first nine coefficients are
\[
 (1,2,22,23,20,4,-2,3,-1).
\]
Thus
\[
 D(x)A(x)=1+2x+22x^2+23x^3+20x^4+4x^5-2x^6+3x^7-x^8.
\]
Negating numerator and denominator gives exactly
\eqref{eq:quartic-gf}. The denominator has constant coefficient $-1$,
so it is a unit in $\mathbb Z[[x]]$ and determines the quotient uniquely.
This proves Theorem~\ref{thm:quartic-main}.

\paragraph{Source indexing.}
The source's generating function has constant coefficient $1$ despite
its offset of one; our convention is explicitly $\sum_{j\ge0}a_{j+1}x^j$.
The separately printed expression
\[
 \frac{47}{9}+14n+\frac{41n^2}{3}+6n^3+n^4
       -\frac29\cos(2\pi n/3)
\]
equals $a_{n+2}$ for positive integer $n$. This follows by substituting
$n+2$ into $Q$ and using the cosine's residue-class values $1,-1/2,-1/2$.
At $n=1$ it gives $40$, whereas $a_1=1$. It must not be interpreted
as a correct same-index formula. The raw recurrence cutoff ``$n>2$''
also fails to specify valid previous positive indices, and even the
well-defined case $n=9$ fails: its right side is $5233$ instead of
$a_9=5232$. The corrected unbounded cutoff is $n\ge10$.
The exact generating-function assertion is the selected resolution;
these index clarifications are not additional problem counts.

\paragraph{Formal scope and prior work.}
Eight modules define the finite reciprocal sum and harmonic mean,
prove the central moments and rational bounds, check the polynomial
certificates, identify the floor at every positive index, and establish
the full formal generating function. The endpoint
\path{QuarticHarmonic.full_resolution} includes the literal coefficient
identification and the exact source numerator--denominator identity.
All principal dependencies are the three standard Lean logical axioms;
there is no admitted inequality, asymptotic hypothesis, or floating-point
calculation in the proof.

The nearby cubic problem A227012 and a general theorem for cubic block
functions already have an exact proof and Lean development in
\path{rajdeepbaror/kimberling-a227012-lean} \cite{prior-cubic-harmonic}.
We inspected that work and excluded the cubic problem from this batch.
Its hypotheses do not cover the present fourth-power blocks.
Moment identities and symmetric reciprocal bounds are elementary
classical tools. Existing jOEIS and LODA code already encodes the proposed
recurrence; the contribution here is its complete proof from the literal
harmonic mean. The bounded exact and semantic prior-work screen is
recorded separately and does not establish worldwide priority.
